% Prefill/decode disaggregation material removed from the paper on 2026-09-23
% for a follow-up paper. Not \input anywhere. Lean proofs (PDDisaggregation.lean,
% capacity only; DecodeScaling.lean, the decode sojourn at equal throughput) and
% simulator checks (validation/src/sim/pd.py, the FIFO tandem without batching;
% validation/src/sim/pd_batching.py, the step engines of serQ #208) remain in the
% repository. Issue #28 (2026-10-01): the throughput no-gain of prop:pd is not
% a latency no-gain; see the "Win condition and latency" and "Limitations"
% paragraphs below.

\begin{proof}[Proof of Proposition~\ref{prop:pd}]
(i) Let $m=\min(N_P/s_P,N_D/s_D)$. Then $ms_P\le N_P$ and $ms_D\le N_D$.
Adding, $m(s_P+s_D)\le N_P+N_D=N$, i.e.\ $m\le N/(s_P+s_D)$.

(ii) At $N_P=Ns_P/(s_P+s_D)$ and $N_D=Ns_D/(s_P+s_D)$ both arguments of
the minimum equal $N/(s_P+s_D)$, so the bound in (i) is attained and no
split does better.

(iii) A number is below a minimum of four terms exactly when it is below
each of them. With $g_P=g_D=1$ and $I=0$ the compute term is
$N/(s_P+s_D)=\mu_A$, and the strict inequality cannot hold, which is (i)
again.
\end{proof}


% ===========================================================================
\section{Prefill/Decode Disaggregation}\label{app:pd}
% ===========================================================================

This appendix collects our results on prefill/decode (PD)
disaggregation. They use the model of \S\ref{sec:model} and are proved
in Appendix~\ref{app:proofs}. A fuller treatment, including batching
inside each pool, is left to follow-up work.

\subsection{Static PD splits}\label{sec:pd}

\paragraph{The claim.} ThunderAgent \citep[App.~A.2, Fig.~7]{thunderagent}
reports that static PD splits of 1:3, 2:2 and 3:1 reach 87.1, 160.1 and
183.2 step/min, against 214.9 for aggregated serving. It attributes the
gap to a smaller effective HBM pool for prefill (``PD disaggregation
exacerbates thrashing'').

\paragraph{The model.} We treat each pool as a capacity and ask when a
split beats aggregation.

\begin{proposition}[PD disaggregation: no-gain theorem and exact win condition]\label{prop:pd}
Let $N$ identical devices serve requests with prefill work $s_P>0$ and
decode work $s_D>0$ (device-seconds). Ideal aggregation has capacity
$\mu_A=N/(s_P+s_D)$; a static split $N_P+N_D=N$ has capacity
$\mu_{PD}=\min(N_P/s_P,\,N_D/s_D)$.
\begin{enumerate}[nosep,leftmargin=1.6em,label=(\roman*)]
\item $\mu_{PD}\le\mu_A$ for every split.
\item Equality holds at the rate-matched split $N_P=Ns_P/(s_P+s_D)$,
      which is therefore optimal.
\item Let aggregation pay an interference overhead~$I$, let dedicated pools
      realise specialisation gains $g_P,g_D$, and let PD be further capped
      by KV-transfer capacity $B_{\mathrm{net}}/\E[K]$ and memory caps
      $\mu_{P,\mathrm{mem}},\mu_{D,\mathrm{mem}}$. Then
      \begin{multline*}
        \mu_{PD} > \mu_A \iff
        \frac{N}{s_P+s_D+I} < \\
        \min\Bigl[\tfrac{N}{s_P/g_P+s_D/g_D},\;
                  \tfrac{B_{\mathrm{net}}}{\E[K]},\;
                  \mu_{P,\mathrm{mem}},\;\mu_{D,\mathrm{mem}}\Bigr],
      \end{multline*}
      i.e.\ PD wins iff \emph{each} of its four bottlenecks exceeds the
      aggregated capacity. With $g_P=g_D=1$ and $I=0$ the compute term
      equals $\mu_A$, recovering~(i).
\end{enumerate}
\provedby{\leanref{pd\_le\_agg}, \leanref{pd\_eq\_agg\_at\_rate\_match}, \leanref{rate\_match\_optimal}, \leanref{pd\_beats\_agg\_iff}, \leanref{pdCompute\_eq\_agg\_of\_no\_gain}}
\end{proposition}


\paragraph{Reading the claim.} Proposition~\ref{prop:pd}(i)--(ii) says the
reported result is what the model predicts under the paper's conditions:
same hardware, no phase specialisation ($g_P=g_D=1$), and a memory-bound
prefill pool (a smaller HBM pool thrashes sooner, i.e.\
$\mu_{P,\mathrm{mem}}$ binds). Nothing specific to agentic workloads is
needed; the bound is a rate-matching identity. Proposition~\ref{prop:pd}(iii)
gives the condition under which the conclusion flips. PD is deployed in
practice \citep{vllm-agentx,mooncake-vllm} with long-context prefill
parallelism ($g_P>1$), wide expert-parallel decode ($g_D>1$), measurable
prefill/decode interference ($I>0$), a fast interconnect
($B_{\mathrm{net}}/\E[K]$ not binding), and a shared KV tier that removes
$\mu_{P,\mathrm{mem}}$ as a constraint. That is the regime the condition
identifies. The supported statement is regime-dependent: static PD
without specialisation and with a memory-bound prefill pool cannot beat
aggregation. The generalisation that PD does not help agentic serving is
not supported.

A second issue is the request shape. For turn $t\ge2$ a request is
``$10^5$ cached tokens $+\,500$ new tokens''. Routing it through the
prefill pool and shipping the whole KV to a decode node is often the
wrong path. Static PD is therefore the wrong abstraction for multi-turn
programs. The routing comparison of \S\ref{sec:routing} makes the
per-turn alternative precise.

\subsection{Per-turn routing between pools}\label{sec:pd-append}

Applied to PD pools, the comparison~\eqref{eq:lookahead} decides where a follow-up turn
is prefilled. Let $W_P,S_P$ be the prefill pool's wait and append service,
$T_{\mathrm{KV}}$ the transfer of the program's KV to the decode node,
$W_D$ the decode wait, $S_D$ the append service on the decode node, and
$I$ the interference it inflicts on co-located decodes.

The decode-local path costs $W_D+S_D+I$ and the prefill-pool path costs
$W_P+S_P+T_{\mathrm{KV}}+W_D$. The decode wait cancels, so decode-local
processing of turn $t\ge2$ is cheaper iff
\begin{equation}\label{eq:append}
  S_D+I < W_P+S_P+T_{\mathrm{KV}} .
\end{equation}
\provedby{\leanref{append\_prefill\_rule}}


Rule~\eqref{eq:append} is the append-prefill rule of PPD
\citep{ppd} written as a cost comparison. Together with
Proposition~\ref{prop:pd} it says that the useful form of PD for agentic
serving is \emph{dynamic}: first-turn prefills go to the prefill pool, and
follow-up turns are routed per turn by comparing the two sides of
\eqref{eq:append}. PPD-style routing is, like session
affinity and KV-aware routing, a one-step approximation
of~\eqref{eq:mdp}.

\subsection{Planned experiment E5: PD win condition}\label{sec:exp-pd}

\emph{Hypothesis} (Prop.~\ref{prop:pd}): the inequality in
Prop.~\ref{prop:pd}(iii), evaluated from independently measured
parameters, predicts which of aggregated and rate-matched PD serving has
higher throughput. We measure $s_P, s_D, I, g_P, g_D$,
$B_{\mathrm{net}}/\E[K]$ and the memory caps, evaluate the condition, and
run both deployments. We repeat with a shared KV tier and with per-turn
append-prefill routing (dynamic PD).

\begin{table}[t]
\caption{E5: predicted vs.\ observed winner. $\mu$ in requests/s.}
\label{tab:pd}
\centering\small
\begin{tabular}{@{}lccc@{}}
\toprule
Configuration & $\mu_A$ & $\min[\cdot]$ & Pred.\ / Obs. \\
\midrule
Static PD, local KV      & \tbd & \tbd & \tbd{} / \tbd \\
Static PD, shared tier   & \tbd & \tbd & \tbd{} / \tbd \\
Dynamic PD (append)      & \tbd & \tbd & \tbd{} / \tbd \\
\bottomrule
\end{tabular}
\end{table}

E5 uses the testbed and protocol of \S\ref{sec:exp}. Its decision
agreement will be scored as in \S\ref{sec:exp-scorecard}.

\subsection{Uncalibrated simulation}\label{sec:sim-pd}

The simulator of \S\ref{sec:sim} also models aggregated and
disaggregated pools. As in \S\ref{sec:sim}, every number below is a
property of the simulated model under a synthetic workload, not a
measurement.

% Generated by `uv run python -m validation.paper_tables` in
% validation/. Do not edit by hand.
\begin{table}[htbp]
\caption{Scalar PD capacities under the assumptions of Prop.~\ref{prop:pd} (in-model): $N=32$, $s_P=s_D=1$, $g_P=2$. Capacities in requests per second.}
\label{tab:sim-pd-inmodel}
\centering\small\setlength{\tabcolsep}{3pt}
\begin{tabular}{@{}l@{\hspace{4pt}}lcc@{}}
\toprule
Quantity & Result & Theory & Simulated \\
\midrule
$\mu_A$, $I=0.5$ & Prop.~\ref{prop:pd} & 12.80 & 12.81 \\
$\mu_{PD}$, $N_P=11$ & Prop.~\ref{prop:pd} & 21.00 & 21.05 \\
\bottomrule
\end{tabular}
\end{table}


\paragraph{In-model.} Table~\ref{tab:sim-pd-inmodel} shows that the
saturated throughput of the aggregated pool and of a disaggregated
tandem matches the scalar capacities of Proposition~\ref{prop:pd}. This
confirms the simulator, not the proposition.

% Generated by `cargo run --release --example paper_tables` in
% libqueuingsim/. Do not edit by hand.
\begin{table}[htbp]
\caption{Latency at equal capacity: $N=8$, $s_P=1$, $s_D=3$, no gains or interference, rate-matched split $N_P=2$. Both configurations have capacity 2 requests/s. Mean end-to-end latency (s).}
\label{tab:sim-pd}
\centering\small\setlength{\tabcolsep}{3pt}
\begin{tabular}{@{}ccc@{}}
\toprule
$\lambda$ & Aggregated & PD ($N_P=2$) \\
\midrule
1.0 & $4.06 \pm 0.02$ & $4.44 \pm 0.03$ \\
1.6 & $4.95 \pm 0.13$ & $7.16 \pm 0.38$ \\
1.9 & $11.62 \pm 3.47$ & $23.00 \pm 5.71$ \\
\bottomrule
\end{tabular}
\end{table}


\paragraph{Win condition and latency.} On a grid of interference $I$,
prefill gain $g_P$ and link capacity $B_{\mathrm{net}}/\E[K]$ ($N=16$),
we compare the simulated winner (aggregation, or the best integer split)
with the condition of Proposition~\ref{prop:pd}(iii). Cells within
3\,\% of the boundary are skipped. The two agree in \simPdAgree{} of
\simPdCells{} decisive cells. The scalar model says nothing about
latency. Table~\ref{tab:sim-pd} shows that at the rate-matched split,
where the capacities are equal, disaggregated serving has a higher mean
latency, and the gap grows with load. A single pool of $N$ servers
queues less than a tandem of two smaller pools, a classical pooling
effect. This is a result about FIFO pools without batching: each request
holds one device for its whole prefill or decode, so the tandem can only
lose. It is the capacity and pooling baseline, not a statement about an
engine. In the simulated step engines of serQ \#208 (vLLM's rules,
exclusive prefill steps, a free transfer, memory that never binds) the
decodes of a colocated device advance only in its decode steps, while a
dedicated decode device steps every iteration and carries more decodes
per step; at the same output throughput the split shortens the time per
output token and lengthens the time to first token. This is a property
of the simulated programs, not a measurement. The processor-sharing
idealisation behind it is exact: scaling the arrival rate and the decode
capacity by one factor keeps the number decoding and divides the sojourn
by that factor.\provedby{\leanref{psMeanNumber\_scale}, \leanref{decode\_sojourn\_scale}, \leanref{dedicated\_sojourn}}
A comparison of PD and aggregation at
equal capacity therefore has to report, separately, the pooling loss of
the tandem, the gain in time per output token, and the cost in time to
first token and transfer.

\paragraph{Limitations.} The PD model treats capacity as a scalar
bottleneck and ignores batching dynamics inside each pool; its no-gain
statement is about requests per second, not about the time a request
spends in the system at a given throughput. The simulator of this
section does not model memory caps on PD pools.

\paragraph{Metrics.} The follow-up reports the following, each defined
once. \emph{TTFT} is measured to the token the client sees. In the model
of \S\ref{sec:pd-append} the prefill produces the first token and the
client receives it after the hand-over, so the transfer and the decode
instance's admission are inside the TTFT (the lecture example records
this as \texttt{ttft\_client}, next to the prefill instance's own
\texttt{ttft}); in an engine that drops the prefill instance's token and
recomputes it, as llm-d does, the first token is the decode instance's
and its first step is inside the TTFT as well. \emph{TPOT} of a request
is the first-to-last span over its gaps,
$(t_{\mathrm{last}}-t_{\mathrm{first}})/(o-1)$ for $o\ge2$ output
tokens; the \emph{request-weighted} mean averages this over requests and
the \emph{token-weighted} mean divides all decode time by all gaps, and
the two differ when short answers wait behind prefills. \emph{Response}
is the time to the last token, TTFT plus the first-to-last span, so no
step is counted twice. \emph{Request
throughput} (requests/s) and \emph{output throughput} (tokens/s) are
measured in a stable load region, where the run keeps up with its
arrivals; the \emph{saturated capacity} is measured with a closed
population and is the quantity of Proposition~\ref{prop:pd}; \emph{SLO
goodput} counts the requests that meet a TTFT and TPOT target. No
steady-state latency is claimed at the rate-matched capacity boundary
itself.

